% ============================================================
%  CAPITOLO 5 --- PROGETTAZIONE DELL'APPLICAZIONE
% ============================================================
\chapter{Progettazione dell'applicazione}

\section{Architettura del sistema}

L'applicazione segue un'architettura \textbf{client-server a tre livelli}:

\begin{enumerate}
  \item \textbf{Presentation layer} --- interfaccia web accessibile tramite browser
  \item \textbf{Application layer} --- server web con logica di business
  \item \textbf{Data layer} --- DBMS PostgreSQL
\end{enumerate}

\begin{figure}[H]
  \centering
  % \includegraphics[width=0.75\textwidth]{immagini/architettura.png}
  \fbox{\parbox{0.75\textwidth}{\centering\vspace{2.5cm}
    \textit{[Inserire diagramma architettura]}\\
    Browser $\longrightarrow$ Flask/Python $\longrightarrow$ PostgreSQL
  \vspace{2.5cm}}}
  \caption{Architettura a tre livelli dell'applicazione}
  \label{fig:architettura}
\end{figure}

\section{Tecnologie utilizzate}

\begin{table}[H]
  \centering
  \caption{Stack tecnologico}
  \label{tab:stack}
  \rowcolors{2}{rowalt}{white}
  \begin{tabularx}{\textwidth}{L{3cm} L{3.5cm} L{7cm}}
    \toprule
    \rowcolor{primary!15}
    \textbf{Livello} & \textbf{Tecnologia} & \textbf{Ruolo} \\
    \midrule
    Database     & PostgreSQL 15   & DBMS relazionale, gestione dati \\
    Backend      & Python 3.11     & Linguaggio applicativo principale \\
    Framework    & Flask 3.x       & Routing HTTP, gestione sessioni, template \\
    ORM          & psycopg2        & Connettore Python--PostgreSQL \\
    Frontend     & HTML5 + CSS3    & Interfaccia utente \\
    Stile        & Bootstrap 5     & Componenti UI responsivi \\
    \bottomrule
  \end{tabularx}
\end{table}

\section{Struttura del progetto}

\begin{verbatim}
stabilimento/
|---- app.py                 # Entry point Flask
|---- config.py              # Configurazione DB e parametri
|---- db/
|   |---- schema.sql         # CREATE TABLE + vincoli
|   `---- seed.sql           # Dati di esempio
|---- routes/
|   |---- auth.py            # Login/logout
|   |---- clienti.py         # CRUD clienti
|   |---- prenotazioni.py    # Prenotazioni e abbonamenti
|   |---- dipendenti.py      # CRUD dipendenti e turni
|   `---- magazzino.py       # Prodotti, fornitori e ordini
|---- templates/
|   |---- base.html          # Layout comune
|   |---- login.html
|   |---- dashboard.html
|   `---- ...                # Una template per sezione
`---- static/
    |---- css/
    `---- js/
\end{verbatim}

\section{Funzionalit\`{a} implementate}

\subsection{Autenticazione}
Il proprietario accede al sistema tramite email e password. La password
viene conservata come hash (bcrypt). La sessione viene mantenuta tramite
cookie firmati di Flask.

\subsection{Gestione clienti}
\begin{itemize}
  \item Inserimento nuovo cliente con validazione email univoca
  \item Modifica e cancellazione (con controllo su prenotazioni esistenti)
  \item Ricerca per nome/cognome
  \item Visualizzazione storico prenotazioni del cliente
\end{itemize}

\subsection{Prenotazioni e abbonamenti}
\begin{itemize}
  \item Selezione cliente e ombrellone con verifica disponibilit\`{a} in tempo reale
  \item Calcolo automatico del prezzo in base a durata e posizione
  \item Inserimento prenotazione con controllo di sovrapposizione date
  \item Gestione degli abbonamenti stagionali
\end{itemize}

\subsection{Gestione dipendenti e turni}
\begin{itemize}
  \item Anagrafica dipendenti
  \item Inserimento e modifica turni con controllo sovrapposizioni
  \item Vista calendario: visualizzazione dipendenti per giornata
\end{itemize}

\subsection{Magazzino e ordini}
\begin{itemize}
  \item Lista prodotti con indicazione delle mancanze (sotto soglia)
  \item Inserimento ordine verso fornitore (multi-prodotto)
  \item Aggiornamento automatico giacenze al momento della conferma consegna
\end{itemize}

% -- 5.5 Screenshot ------------------------------------------
\section{Screenshot dell'interfaccia}

\begin{figure}[H]
  \centering
  % \includegraphics[width=0.9\textwidth]{immagini/screen_login.png}
  \fbox{\parbox{0.9\textwidth}{\centering\vspace{3cm}
    \textit{[Screenshot: pagina di login]}
  \vspace{3cm}}}
  \caption{Pagina di login}
  \label{fig:screen-login}
\end{figure}

\begin{figure}[H]
  \centering
  % \includegraphics[width=0.9\textwidth]{immagini/screen_dashboard.png}
  \fbox{\parbox{0.9\textwidth}{\centering\vspace{3cm}
    \textit{[Screenshot: dashboard principale]}
  \vspace{3cm}}}
  \caption{Dashboard --- riepilogo giornaliero}
  \label{fig:screen-dashboard}
\end{figure}

\begin{figure}[H]
  \centering
  % \includegraphics[width=0.9\textwidth]{immagini/screen_prenotazioni.png}
  \fbox{\parbox{0.9\textwidth}{\centering\vspace{3cm}
    \textit{[Screenshot: gestione prenotazioni]}
  \vspace{3cm}}}
  \caption{Pagina gestione prenotazioni con mappa ombrelloni}
  \label{fig:screen-prenotazioni}
\end{figure}

\begin{figure}[H]
  \centering
  % \includegraphics[width=0.9\textwidth]{immagini/screen_turni.png}
  \fbox{\parbox{0.9\textwidth}{\centering\vspace{3cm}
    \textit{[Screenshot: gestione turni]}
  \vspace{3cm}}}
  \caption{Vista turni dipendenti per giornata}
  \label{fig:screen-turni}
\end{figure}

\begin{figure}[H]
  \centering
  % \includegraphics[width=0.9\textwidth]{immagini/screen_magazzino.png}
  \fbox{\parbox{0.9\textwidth}{\centering\vspace{3cm}
    \textit{[Screenshot: magazzino]}
  \vspace{3cm}}}
  \caption{Pagina magazzino con evidenza prodotti sotto soglia}
  \label{fig:screen-magazzino}
\end{figure}
